\documentclass[a4paper,11pt,oneside]{article}
\usepackage{tabularx}
\usepackage{fancybox}
\usepackage{amsmath}
\usepackage[top=1.5in,bottom=1in,right=1in,left=1in,headheight=65pt]{geometry}
\usepackage{graphicx}
\usepackage{fancyhdr}

%Gummi|065|=)
\begin{document}
\pagestyle{fancy}

\fancyhead[L]{Stevens UPEPC\\12 April 2014}
\fancyhead[C]{{\LARGE\textbf{Time Out}}}
\fancyhead[R]{\includegraphics[height=32pt]{../stevenslogo.eps}\ \ \ \includegraphics[height=32pt]{../upelogo.png}}


\section{Problem Statement}

Steven's digital clock seems to be malfunctioning. When he checked the time at various
points throughout the day, he noticed that some of the LCD segments are not working.
\\\\
The displays he has seen so far have been:
\\\\
\begin{tabularx}{\textwidth}{X X X X}
\includegraphics[height=40pt]{./assets/testclock_1.eps}&
\includegraphics[height=40pt]{./assets/testclock_2.eps}&
\includegraphics[height=40pt]{./assets/testclock_3.eps}&
\includegraphics[height=40pt]{./assets/clock_4.eps}
\end{tabularx}
\\\\
Based on his previous observations, can Steven accurately determine the current time,
shown by last clock display? More generally, given the previous observations and the
current clock display, your job is to determine the current time, unless it is not
possible to do with certainty.

\section{Input}
The first line of input contains a single integer $\boldsymbol{P}$,
$(\boldsymbol{1} \le \boldsymbol{P} \le \boldsymbol{1000})$, which is the
number of data sets that follow. Each data set should be processed identically
and independently.
\\\\
Each data set begins with a single line that contains $\boldsymbol{K}$, the data
set number, followed by $\boldsymbol{N}$,
$(\boldsymbol{0} \le \boldsymbol{N} \le \boldsymbol{10})$ which is the number
of clock displays that Steven has seen previously. Note that Steven does not
remember the clock displays in order and he may have seen them any time in the past
12 hours. Also note that Steven's clock uses a 12-hour format.
\\\\
The next $\boldsymbol{N}$ lines  contain the clock displays that Steven has seen.
The clock displays are given in the format $(d_1\ d_2\ d_3\ d_4)$ where $d_0$ is
the tens digit of the hour and $d_4$ is the ones digit of the minutes. Being as
the digits in the display may not be complete digits, each will be given as
a binary representation of the LCD segments that are active (see below). The last
line of each data set is of the same $(d_1\ d_2\ d_3\ d_4)$ format and indicates
the current clock time, $T$.

\subsection{Digit Representation}
\begin{tabularx}{\textwidth}{X r}
\parbox[m]{10cm}{Digits are binary-encoded. Being as there are 7 LCD segments that
make up each digit, they will be encoded as an integer value such that the $i$th segment
is active in a display $D$ iff the $i$th least significant bit is set in $D$.
\\\\
For example, the digit ``5'' would be encoded as:
\begin{center}
$2^0 + 2^1 + 2^2 + 2^5 + 2^6 = 103_{10}$
\end{center}
} & \parbox[m]{5cm}{\includegraphics[height=120pt]{./assets/digit_binary.eps}} \\
\end{tabularx}

\section{Output}
For each data set there is a single line of output. The single line of output
consists of the data set number $\boldsymbol{K}$, followed by a single space
followed by the real value of the current clock time $\boldsymbol{T}$. The time you output
should be in regular human-readable 12-hour clock format, i.e ``\texttt{11:59}''
or ``\texttt{2:30}''. If it is impossible to determine with 100\% certainty
what the current time is, your program should output the string ``\texttt{NO SOLUTION}''
instead (without quotes).

\section{Test Data}
\emph{TODO}\\\\
\begin{tabularx}{\textwidth}{|X|X|}
	\hline
	Input & Output \\ \hline
	\parbox[t]{5cm}{
	\texttt{3\\
			1\\
			2\\
			3\\
	}} & \parbox[t]{5cm}{
	\texttt{1 ???\\
			2 ???\\
			3 ???\\
	}}\\
	\hline
\end{tabularx}

\end{document}








































